\documentclass[11pt,a4paper]{article}
\usepackage{fontspec}
\setmainfont{Noto Serif}
\setsansfont{Noto Sans}
\usepackage{babel}
\babelprovide[main,import]{russian}
\usepackage{amsmath,amssymb,mathtools,booktabs,geometry,fancyhdr,microtype,xcolor,hyperref}
\geometry{left=21mm,right=21mm,top=21mm,bottom=21mm}
\definecolor{ink}{HTML}{173D57}
\hypersetup{colorlinks=true,urlcolor=ink,linkcolor=ink,pdftitle={BGS-F. Арифметика точных спектральных совпадений}}
\pagestyle{fancy}\fancyhf{}\lhead{\small BGS-F}\rhead{\small Арифметика точных спектральных совпадений}\cfoot{\thepage}
\setlength{\headheight}{15pt}\setlength{\parindent}{1em}\setlength{\parskip}{2.5pt}
\setlength{\emergencystretch}{3em}\tolerance=3000
\newcommand{\Tr}{\operatorname{Tr}}
\newcommand{\EE}{\mathbb E}
\newcommand{\R}{\mathbb R}
\newcommand{\Z}{\mathbb Z}
\newcommand{\C}{\mathbb C}
\newcommand{\si}{\mathrm{Sine}_2}
\newcommand{\eps}{\varepsilon}
\newcommand{\sinc}{\operatorname{sinc}}
\newcommand{\sect}[1]{\section*{\sffamily\color{ink}#1}}
\newcommand{\subt}[1]{\par\medskip\noindent\textbf{#1.}\quad}
\newcommand{\T}{\mathbb T}
\newcommand{\dist}{\operatorname{dist}}
\newcommand{\Ad}{\operatorname{Ad}}
\newcommand{\zp}{\mathbb Z_7}
\newcommand{\QQ}{\mathbb Q}
\newcommand{\kk}{\overline{\mathbb Q}}
\newcommand{\ord}{\operatorname{ord}}
\newcommand{\diag}{\operatorname{diag}}
\begin{document}
\noindent{\small\sffamily\color{ink} МАТЕМАТИЧЕСКАЯ ФИЗИКА\hfill 20 НОЯБРЯ 2026}
\begin{center}{\LARGE\sffamily Кратности при фиксированном кике}\\[4pt]
{\large Алгебраические разности фаз и расщепление на степенях семи}
\end{center}
Пусть \(U_N(t)=e^{itA_N}M_N\), где \(A_N,M_N\) взяты из исходной задачи. Целевая последовательность соответствует \(t=N/10^{80}\). Обозначим \(o_N=\ord(S\bmod N)\), а через \(s_N\) --- число простых собственных значений целевого \(U_N\).

\textbf{Алгебраический перенос.} При любом ненулевом вещественном алгебраическом \(t\) множество кратностей \(U_N(t)\) совпадает с общим случаем аналитического семейства. Если \(\lambda,\mu\) --- собственные значения \(U_N(t)\) и \(\lambda/\mu\) алгебраично, то \((\lambda/\mu)^{o_N}=1\). В частности, для \(U_N(t)^{o_N}\) отношения \emph{различных} собственных значений трансцендентны. Перенос общего случая относится также к рангу \(\Ad(U_N(t)^m)-zI\) для каждого \(m\in\Z\) и каждого алгебраического \(z\).

\textbf{Нерасщеплённые модули.} При \(N=3^k\) или \(N=5^k\) для каждого \(k\ge1\) спектр исходного \(U_N\) прост. Доказательство справедливо и для остальных простых \(p\) с указанным в разделе 2 условием порядка. Также \(U_N^m\) имеет простой спектр при \(\gcd(m,o_N)=1\). На этих размерностях точное среднее \(|\Tr U_N^m|^2/N\) при \(m=1,\ldots,T\) стремится к \(1\) при \(T\to\infty\) для каждого фиксированного \(N\). Случаи \(p=3,5\) уже были установлены в предыдущей части о совпадениях. Здесь они включены как следствие общего критерия.

\textbf{Расщеплённые модули.} Для \(N=7^k\) у невозмущённой карты есть \(18N/49\) двойных уровней с примитивным \(7\)-адическим индексом (\(k\ge2\)). Доля таких уровней, которые расщепляются в первом порядке, равна \(83/96+O(N^{-1/2})\). В кластерах каждого фиксированного положительного \(7\)-адического порядка получены ещё четыре простые ветви. Алгебраический перенос даёт для исходного фиксированного кика
\[
 \liminf_{k\to\infty}\frac{s_{7^k}}{7^k}\ge\frac{345}{392},\qquad
 \limsup_{k\to\infty}\lim_{T\to\infty}\frac1{T7^k}
              \sum_{m=1}^T|\Tr U_{7^k}^m|^2\le\frac{64549}{57624}.
\]
Вторая константа примерно \(1.1202\). Для сравнения, предельное число для линейной модели при том же порядке пределов равно \(8/3\).

\textbf{Обычный след при \(m=3N/7\).} Для \(N=7^k\), \(k\ge3\), получена точная кривизна в точке линейного возврата. Если \(I_N(u)=N^{-2}|\Tr(e^{iNu A_N}M_N)^{3N/7}|^2\), то
\[
               I_N''(0)=-\frac{34112649}{14000000}\,N^3 .
\]
Здесь \(u\) --- физическая сила. Формула относится к производной при нуле. Для предельного парного теста при \(u=10^{-80}\) по-прежнему требуется количественный контроль ненулевых разностей фаз масштаба \(1/N\) и обычных следов при \(m\asymp N\).
\newpage
\sect{1. Перенос ранга на фиксированный алгебраический параметр}
Пусть \(A\) --- эрмитова алгебраическая матрица, \(M\) --- алгебраическая унитарная матрица. Алгебраичность означает, что все элементы лежат в \(\kk\). Все собственные числа \(a_j\) матрицы \(A\) алгебраичны. Если в списке \(a_j\) оставить различные, спектральный проектор для \(a_j\) выражается как \(\prod_{l\ne j}(A-a_l I)/(a_j-a_l)\), поэтому также алгебраичен. Отсюда элементы \(U(t)=e^{itA}M\) и \(U(t)^{-1}\) принадлежат кольцу конечных сумм
\[
 \mathcal R=\left\{\sum_{\alpha}c_\alpha e^{it\alpha} :
              c_\alpha\in\kk,\ \alpha\in\kk\cap\R\right\}.
\]
После группировки равных частот коэффициенты определены однозначно. По теореме Линдемана--Вейерштрасса вычисление в любой точке \(t_0\in\kk\cap\R\), \(t_0\ne0\), является инъективным на \(\mathcal R\): показатели \(it_0\alpha\) различны и алгебраичны. Применяя инъективность к каждому минору, получаем, что \(\operatorname{rank}B(t_0)=\max_{t\in\C}\operatorname{rank}B(t)\) для любой матрицы над \(\mathcal R\). Каждый ранг достигается в общем случае и на \(\R\), так как ненулевой аналитический минор имеет дискретное множество нулей. Для исходной модели \(M_N\) алгебраична по явной формуле, а \(A_N\) алгебраична по рациональным комплексным коэффициентам ряда Фурье. Значение \(t=N/10^{80}\) удовлетворяет условию.

\subt{Кратности}
Возьмём матрицу \(H_{ab}(t)=\Tr U(t)^{a+b}\), \(0\le a,b<N\). Если различные собственные числа равны \(\lambda_j\) и их алгебраические кратности равны \(d_j\), то
\[
                       H=V\,\diag(d_j)\,V^T,\qquad V_{aj}=\lambda_j^a.
\]
Формула верна и при жордановых блоках. Матрица Вандермонда имеет полный столбцовый ранг, поэтому ранг \(H\) есть точное число различных корней. Там, где этот ранг максимален, кратности локально постоянны как мультимножество: малые непересекающиеся окружности вокруг корней сохраняют их полные кратности, а число различных корней уже максимально. Исключительное множество дискретно, и его дополнение в \(\C\) связно. Следовательно, все точки максимального ранга имеют одну и ту же разбивку на кратности. Любой простой корень при произвольном параметре сохраняется в малом круге при переходе в общий случай. В частности, общее число простых уровней не меньше, чем в любой отдельно рассмотренной точке.

\subt{Алгебраические отношения собственных чисел}
Элементы \(\Ad(U(t)):\ X\mapsto U(t)XU(t)^{-1}\) также принадлежат \(\mathcal R\). Для любого \(z\in\kk\) точно
\[
       \dim\ker\big(\Ad(U(t_0))-zI\big)
               \le \dim\ker\big(\Ad(M)-zI\big).
\]
Правая часть нулевая, если \(z^{o_N}\ne1\), поскольку \(\Ad(M_N)\) переставляет базис \(T_N(v)\) по \(S\) и имеет порядок \(o_N\). Для вещественного \(t\) собственные числа \(\Ad(U(t))\) равны \(\lambda_i/\lambda_j\). Тот же довод работает для любой степени этой перестановки. В частности, при \(m=o_N\) её алгебраический спектр состоит из одной единицы. Ранее в части о жёсткости Вейля получено более общее утверждение о специализации полиномиальных отношений. В частности, любой алгебраический сбалансированный моном \(\prod_{h=1}^r(\lambda_{i_h}/\lambda_{j_h})\) обязан быть корнем из единицы порядка, делящего \(o_N\), по применению этого доказательства к \(\Ad(U(t))^{\otimes r}\). Источник используемой трансцендентности --- \href{https://encyclopediaofmath.org/wiki/Lindemann_theorem}{Encyclopedia of Mathematics, Lindemann theorem}.
\newpage
\sect{2. Простой спектр на нерасщеплённых модулях}
Пусть \(p\) нечётно, \(2\) не квадрат по модулю \(p\), \(S\) имеет порядок \(p+1\) по модулю \(p\), а \(S^{p+1}\not\equiv I\bmod p^2\). Отождествим действие \(S\) с умножением на \(\gamma=3+2\sqrt2\) в неразветвлённом кольце \(\mathbb Z_p[\sqrt2]\), используя координату \(v_1+\sqrt2 v_2\). Имеем
\(\gamma^{p+1}=1+p b\) с обратимым \(b\). Бином Ньютона при нечётном \(p\) даёт
\(\ord(\gamma\bmod p^a)=(p+1)p^{a-1}\) для всех \(a\ge1\). Любой вектор аддитивного порядка \(p^a\) в \((\Z/p^k\Z)^2\) представлен вектором, делящимся на \(p^{k-a}\), после деления на который его координата в кольце является единицей. Поэтому такой вектор имеет длину орбиты ровно \((p+1)p^{a-1}\).

\subt{Подсчёт орбит}
Каждая орбита вейлевого базиса даёт по одному инвариантному вектору при \(\Ad M_N\). Их полное число
\[
 1+\sum_{a=1}^k\frac{p^{2a}-p^{2a-2}}{(p+1)p^{a-1}}
          =1+(p-1)\sum_{a=1}^kp^{a-1}=p^k=N.
\]
Для унитарной матрицы с кратностями \(d_j\) размер коммутанта равен \(\sum_j d_j^2\). Минимальное значение \(N\) эквивалентно простоте. Это доказывает простоту \(M_N\); для \(U_N(t_0)\) применяем раздел 1 с \(z=1\). Если \(\gcd(m,o_N)=1\), то \(M_N^m\) также имеет простой спектр, так как возведение в \(m\)-ю степень инъективно на корнях из единицы порядка, делящего \(o_N\), с точностью до общей фазы. Применяем перенос к \(U_N(t_0)^m\).

\subt{Явные простые}
При \(p=3\) порядок \(S\bmod3\) равен 4 и \(S^2+I=6S\); множитель \(S^2-I\) обратим по модулю \(3\). При \(p=5\) порядок \(S\bmod5\) равен 6 и
\[
                 S^3+I=\begin{pmatrix}100&140\\70&100\end{pmatrix}.
\]
Множитель \(S^3-I\) обратим по модулю \(5\), поэтому условие поднятия также выполняется. Точный перебор простых \(p<1000\) подтверждает условие для 66 простых, включая \(3,5,11,19,37,43,53\).

\subt{Длительное среднее}
Для любой фиксированной унитарной матрицы \(U\) с кратностями \(d_j\) геометрическая сумма по временам даёт
\begin{equation}
  \lim_{T\to\infty}\frac1{TN}\sum_{m=1}^T|\Tr U^m|^2
          =N^{-1}\sum_j d_j^2.                                  \tag{1}
\end{equation}
Если \(U=U_N(t_0)\), то эта сумма равна минимальному значению коммутантной размерности семейства, делённому на \(N\). При простоте она равна \(1\). Здесь \(T\to\infty\) при фиксированном \(N\).
\newpage
\sect{3. Двойные уровни на степенях семи}
Пусть \(N=7^k\), \(k\ge2\), \(p=7\). Берём \(\rho\in\zp\) с \(\rho^2=2\), \(\rho\equiv3\bmod7\). Симплектическая замена на вейлевых метках
\[
     (v_1,v_2)\longmapsto (a,b)=
     \left(v_1+\rho v_2,\frac{v_1-\rho v_2}{-2\rho}\right)
\]
имеет определитель \(1\), диагонализует \(S\) и даёт \(\gamma=3+2\rho\) на первой координате. Сохранение симплектической формы обеспечивает, что замена меток реализуется точным унитарным сопряжением базиса Вейля. Действительно, она сохраняет правило умножения, след, звёздочку и образует \(*\)-автоморфизм полной матричной алгебры, а любой такой автоморфизм внутренний. В этих координатах \(M_N\) с точностью до общей фазы действует как \(\psi(x)\mapsto\psi(\gamma^{-1}x)\), поскольку оба оператора одинаково сопрягают все операторы Вейля.

Порядок \(\gamma\) равен \(o=3\cdot7^{k-1}\). По модулю \(7\) его порядок равен 3, а утверждение для последующих степеней следует из
\[
                    S^3-I=14\begin{pmatrix}7&10\\5&7\end{pmatrix}
\]
тем же биномиальным поднятием. Единицы распадаются на две орбиты, чей первый разряд принадлежит квадратам \(\{1,2,4\}\) или неквадратам \(\{3,5,6\}\). Пусть \(\omega:\mathbb F_7^\times\to\{1,\zeta_3,\zeta_3^2\}\) --- характер, \(\omega(3)=\zeta_3\), и
\[
 \mathcal L(x)=\tfrac16\log(x^6),\qquad
 \chi_{\ell,h}(x)=\omega(x\bmod 7)^h\,\exp\big(2\pi i\ell\,\mathcal L(x)/7^k\big).
\]
В экспоненте используется класс \(\mathcal L(x)\bmod7^k\); \(\ell\) перебирает \(\Z/7^{k-1}\Z\), \(h=0,1,2\). Это все характеры единиц, тривиальные на \(-1\). Действительно, логарифм идентифицирует \(1+7\zp\) с \(7\zp\), а ограничение на группу \(\langle\gamma\rangle\) даёт биекцию с её характерами. Если \(7\nmid \ell\), собственное число соответствующего характера имеет полный \(7\)-адический порядок \(7^{k-1}\) после снятия общей фазы. На остальных орбитах \(x\) этот порядок меньше. Поэтому такое собственное число имеет кратность точно два. Число этих двойных уровней \(3(7^{k-1}-7^{k-2})=18N/49\).

Пусть \(u_+,u_-\) --- векторы, равные \(\chi_{\ell,h}/\sqrt{o}\) на двух классах и нулю вне них. Чётность \(P\psi(x)=\psi(-x)\) меняет \(u_+\) и \(u_-\) местами. При \(U_N'(0)=i A_N M_N\) производные поднятых углов в данном кластере равны собственным числам сжатия \(A_N\) на \(\operatorname{span}(u_+,u_-)\). Это следует дифференцированием локального аналитического эрмитова логарифма \(U_N(t)\), после выбора разреза вне спектра. Сжатие нескалярно в точности тогда, когда хотя бы одно из двух чисел
\[
         \langle u_+,A_Nu_-\rangle,\qquad
         \langle u_+,(A_N-PA_NP)u_+\rangle
\]
ненулевое. Для диагональной разности в сумме Вейля берутся коэффициенты \(c_v-c_{-v}=2i\Im c_v\); для внедиагонального элемента берутся \(c_v\). Пара \(v=(0,\pm7)\) имеет вещественные коэффициенты, а в новых метках обе координаты делятся на \(7\). Она не входит ни в один из этих двух тестов.
\newpage
\sect{4. Точная стационарная фаза}
Для каждой оставшейся метки \(v\) обе координаты \(a,b\) обратимы по модулю \(7\). Её вклад в один из двух элементов сжатия, с точностью до общего множителя \(1/o\), есть коэффициент кика, умноженный на сумму по подходящим классам \(x\bmod7\) точного выражения
\[
  w_h(x)\,\zeta_{7^k}^{\Phi(x)},\qquad
  w_h(x)=\left(\frac{\omega(x+a)}{\omega(x)}\right)^h,\qquad
  \Phi(x)=b(x+a/2)+\ell\big[\mathcal L(x+a)-\mathcal L(x)\big].
\]
В обоих тестах \(x\bmod7\) квадрат. Во внедиагональном тесте \(x+a\bmod7\) неквадрат, в диагональной разности --- квадрат. Вес \(w_h\) постоянен на каждом таком классе. Производная
\[
                 \Phi'(x)=b-\frac{\ell a}{x(x+a)}
\]
показывает, что стационарные точки задаются сравнением \(b x(x+a)=\ell a\). Чтобы увидеть точное исчезновение нестационарных классов, пишем \(x=x_0+7^{\lceil k/2\rceil}y\). По ряду логарифма \(\Phi(x)\equiv\Phi(x_0)+7^{\lceil k/2\rceil}y\Phi'(x_0)\bmod7^k\). Сумма по \(y\) нулевая, если \(\Phi'(x_0)\not\equiv0\bmod7^{\lfloor k/2\rfloor}\).

\subt{Окрестность точного корня}
Пусть \(\Phi'(x_*)=0\) в \(\zp\), \(\nu_7(\Phi''(x_*))=d\), и \(k\ge 3d+2\). На классе \(x=x_*+7^{d+1}y\) разложение имеет вид
\[
 \Phi(x)=\Phi(x_*)+7^{3d+2}\,c\, y^2(1+7 y B(y)),
       \qquad c\in\zp^\times .
\]
Действительно, коэффициент при \(y^2\) имеет указанный порядок, а для степеней \(j\ge3\) ряд логарифма даёт порядок не меньше \(j(d+1)-\nu_7(j)\ge 3d+3\). Эти порядки растут, поэтому \(B\) задаётся сходящимся рядом с целыми \(7\)-адическими коэффициентами. Замена \(y\mapsto y\sqrt{1+7yB(y)}\) биективна по каждому модулю \(7^n\) по последовательному поднятию, поскольку она тождественна вместе с производной по модулю 7. Вклад класса равен
\[
         w_h(x_*)\,7^{2d+1}\,\zeta_{7^k}^{\Phi(x_*)}\,
                     G_{7^{k-3d-2}}(c),                        \tag{2}
\]
где \(G_{7^n}(c)=\sum_{y\bmod7^n} e^{2\pi i c y^2/7^n}\), \(G_1=1\). Рекурсия \(G_{7^n}=7G_{7^{n-2}}\) при \(n\ge2\) следует разбиением \(y=y_0+7^{n-1}y_1\). Для \(n=1\) квадрат модуля равен \(7\) по биективной замене \((y,z)\mapsto(y-z,y+z)\). Таким образом, множитель при \(\zeta_{7^k}^{\Phi(x_*)}\) принадлежит \(\QQ(i,\zeta_3,\zeta_7)\) и ненулевой, включая ненулевой коэффициент метки.

\subt{Независимость фаз}
При \(k-r\ge1\) элементы \(\zeta_{7^k}^j\), \(0\le j<7^r\), образуют базис над \(\QQ(i,\zeta_3,\zeta_{7^{k-r}})\) по степени циклотомического расширения. Поэтому вклад (2), чья фаза \(\Phi(x_*)\bmod7^r\) единственна, не может сократиться. Для невырожденных корней \(d=0\) значение \(x_*\bmod7^r\) находится решением квадратного сравнения. Дискретный критерий использует лишь малые модули. Стандартный метод точной стационарной фазы см. в работе \href{https://doi.org/10.1006/jnth.2002.2790}{B. Fisher, J. Number Theory 96 (2002), 201--224}.
\newpage
\sect{5. Невырожденные классы и конечная проверка}
При \(\ell\bmod7\in\{1,2,3,5,6\}\) вычисления над \(\Z/343\Z\) находят хотя бы один класс фазы, который входит ровно один раз в одном из двух тестов. В этом классе \(\Phi(x_*)\not\equiv0\bmod7\), поэтому корень невырожден: \(2x_*+a\not\equiv0\bmod7\). В самом деле, \(\mathcal L\equiv0\bmod7\), а \(b\) обратима. Любой вырожденный стационарный вклад имеет фазу ноль по модулю 7 и не может уничтожить свидетель. Даже без формулы (2) для такого вклада это следует из того, что все его экспоненты имеют показатель кратный \(7\). Они лежат в подпространстве, порождённом базисными элементами с \(j\equiv0\bmod7\).

Точный алгоритм проверяет каждую метку и каждый обратимый \(x\bmod343\), для которого \(x+a\) также обратим. Вычисляются
\[
  \ell_0=b x(x+a)a^{-1}\bmod343,\qquad
  j=b(x+a/2)+\ell_0\big(\mathcal L(x+a)-\mathcal L(x)\big)\bmod343 .
\]
Для диагональной разности исключаются вещественные коэффициенты, затем ведётся кратность каждого \(j\) отдельно для пары \((\ell_0,\text{тест})\). Для каждого \(\ell_0\) из пяти нужных классов найден \(j\) кратности \(1\) с \(7\nmid j\). Используется
\(\mathcal L(x)\equiv\tfrac16[(x^6-1)-(x^6-1)^2/2]\bmod343\).
Число проверенных \(\ell_0\) равно \(245\). Ниже приведены примеры из таблицы точных сертификатов. Тип \(\mathrm{D}\) означает диагональную разность, тип \(\mathrm{O}\) --- внедиагональное сжатие.
\begin{center}
\begin{tabular}{ccccc}
\toprule
\(\ell_0\) & Тип & Метка \(v\) & \(x\bmod343\) & \(j\bmod343\) \\
\midrule
1 & \(O\) & \((0,-9)\) & 296 & 78\\
2 & \(D\) & \((-1,0)\) & 135 & 142\\
3 & \(D\) & \((0,-8)\) & 312 & 193\\
5 & \(D\) & \((0,-9)\) & 15 & 23\\
6 & \(O\) & \((1,0)\) & 114 & 132\\
\bottomrule
\end{tabular}
\end{center}
При \(k\ge4\) поле для \(r=3\) содержит \(\zeta_7\), и (2) применима к указанному корню. Для каждого примитивного характера с \(\ell\bmod7\) в пяти классах сжатие точно нескалярно, независимо от \(h\).

\subt{Малые модули}
Для \(k=2,3\) в архиве сжатие проверяется дополнительно без стационарной фазы. Каждый коэффициент суммы раскладывается по базису
\(i^a\zeta_3^{\,b}\zeta_{7^k}^{\,j}\) с \(a,b\in\{0,1\}\) и
\(0\le j<6\cdot7^{k-1}\), используя
\(\zeta_3^2=-1-\zeta_3\) и
\(\sum_{q=0}^6\zeta_{7^k}^{q7^{k-1}}=0\).
Проверка целочисленная. Для \(k=2\) расщепляются 12 из 18 двойных уровней, для \(k=3\) --- 114 из 126. Эти конечные исключения не входят в асимптотическое доказательство.
\newpage
\sect{6. Вырожденный класс и доля расщеплений}
Для \(\ell\equiv4\bmod7\) в диагональной разности нет стационарных классов. Во внедиагональном элементе они есть ровно у трёх меток \(v=(0,1),(0,8),(0,-6)\), по одному классу у каждой, с \(2x+a\equiv0\bmod7\). Проверка подстановкой по модулю 7 входит в код. Ставим \(y=x+a/2\). Уравнение принимает вид \(y^2=D\), где
\[
             D=a^2/4+\ell a/b,\qquad
             \ell_{\mathrm c}=-ab/4 .
\]
Для трёх меток значения \(\ell_{\mathrm c}\bmod49\) равны \(11,18,4\). При \(k\ge4\) у всякой метки с \(\ell\not\equiv\ell_{\mathrm c}\bmod49\) вклад равен нулю: \(D\) имеет порядок 1, и квадратного корня по модулю 49 нет. Поэтому при каждом \(\ell\) к сумме может вносить вклад только одна из трёх меток.

Если \(v=\nu_7(D)<\lfloor k/2\rfloor\) нечётен или \(7^{-v}D\) не квадрат по модулю 7, вклад точно нулевой. В оставшемся случае \(v=2d\ge2\) и есть два корня \(y=\pm z\) в \(\zp\), где \(\nu_7(z)=d\). При \(k\ge4d+2\) условие стационарности по модулю \(7^{\lfloor k/2\rfloor}\) требует \(y\equiv z\) или \(y\equiv-z\bmod7^{d+1}\). Это следует из \((y-z)(y+z)\equiv0\bmod7^{2d+1}\) и \(\nu_7(2z)=d\). В обеих окрестностях \(\nu_7(\Phi''(x_*))=d\), так что справедлива (2).

\subt{Различие двух критических фаз}
В координате \(y\) производная равна \(b(y^2-z^2)/(y^2-a^2/4)\). Разворачиваем обратный знаменатель в степенной ряд на \(y\in7\zp\) и берём разность первообразных в \(z\) и \(-z\). Ведущий член равен \(-4c z^3/3\) с \(c=-4b/a^2\) обратимым. Для \(j\ge1\) последующие слагаемые имеют порядок не меньше
\[
 (2j+3)d-\max(\nu_7(2j+1),\nu_7(2j+3))>3d.
\]
Последнее неравенство следует из \(7^{2j}>2j+3\). Первые два критических значения, следовательно, различны по модулю \(7^{3d+1}\). При \(k\ge4d+2\) имеем \(k-(3d+1)\ge1\), так что по независимости в разделе 4 два ненулевых вклада не сокращаются. В этом случае двойной уровень расщепляется для всякого \(h\).

\subt{Подсчёт}
Условно по \(\ell\equiv4\bmod7\) и для одной метки \(\ell-\ell_{\mathrm c}\) равномерно в \(7\zp\); умножение на \(a/b\) сохраняет эту равномерность. Доля с \(\nu_7(D)=2d\) и квадратным ведущим разрядом равна \((1-1/7)7^{-(2d-1)}/2\). Сумма по \(d\ge1\) равна \(1/16\), а три успешных множества не пересекаются по модулю 49. При конечном \(k\) пропускаем разряды с \(\nu_7(D)>2\lfloor(k-2)/4\rfloor\); их доля \(O(7^{-k/2})\). Для всех остальных классификация точна. С учётом пяти невырожденных классов доля расщеплений среди всех примитивных двойных уровней равна
\[
                  \frac{5+3/16}{6}+O(7^{-k/2})
                      = \frac{83}{96}+O(N^{-1/2}) .
\]
Наличие нескалярного сжатия переносится на пару простых уровней при достаточно малом ненулевом \(t\) для каждого фиксированного \(N\). Затем применяется раздел 1 к общему, в частности исходному алгебраическому, параметру.
\newpage
\sect{7. Кластеры более высокого \(7\)-адического порядка}
Пусть \(\ell=7^d\mu\), \(7\nmid\mu\), где \(d\ge1\) фиксировано и \(k\ge2d+6\). Линейный кластер имеет размер \(2(d+1)\). По два вектора на каждом уровне \(r\in\{0,\ldots,d\}\) поддержаны на \(x=7^r X\) с данным квадратным классом \(X\bmod7\). Их фаза есть \(\omega(X)^h\zeta_{7^k}^{\,\ell\mathcal L(X)}\), нормировка равна \((3\cdot7^{k-r-1})^{-1/2}\). Ограничение на характер \(\gamma\) даёт нужное собственное число; на глубинах \(r>d\) его порядок по \(7\) недостижим.

\subt{Какие блоки могут быть ненулевыми}
Для примитивной метки \(a,b\) обратимы. Если обе координаты \(x,x+a\) --- единицы, производная \(\Phi'(x)\equiv b\bmod7\), так что сумма нулевая. Две координаты положительной глубины несовместимы с обратимым \(a\). Для блока глубин \(0,r>0\) пишем \(x=-a+7^r Y\); фаза равна
\[
 -ab/2+7^r bY+7^d\mu\big[\mathcal L(Y)-\mathcal L(-a+7^rY)\big].
\]
После снятия общей константы и \(7^{\min(r,d)}\) производная обратима по модулю 7, если \(r\ne d\). Например, при \(r<d\) остаётся \(b\), а при \(r>d\) --- \(\mu/Y\). Поэтому остаётся только связь \(0,d\) и сопряжённый блок. Длина суммы в \(Y\) равна \(7^{k-r}\), и при факторе \(7^{\min(r,d)}\) она покрывает целое число периодов. При указанных \(k\) стационарное условие раздела 4 применимо.

Две непримитивные метки имеют вид \(a=7a_0,\ b=7b_0\). На глубине 0 нормированная производная после снятия \(7\) равна \(b_0\bmod7\); смешение глубины 0 с большей невозможно по поддержке. Если обе координаты имеют глубину 1, после записи \(x=7X\) получается вейлева фаза с множителем \(7^2\), а логарифмическая с множителем \(7^d\). Если \(d\ne2\), после снятия меньшей степени 7 производная обратима. Здесь логарифмическая производная \(-\mu a_0/[X(X+a_0)]\) также обратима. Между глубиной 1 и \(r\ge2\) пишем \(x=7(-a_0+7^{r-1}Y)\). Слагаемые производной различаются по порядку, если \(r+1\ne d\). Между двумя глубинами не меньше 2 переходов нет, поскольку \(a\) имеет порядок 1. Следовательно, при \(d=1\) сжатие целиком имеет вид
\(\left(\begin{smallmatrix}0&C\\C^*&0\end{smallmatrix}\right)\) на глубинах \(0,d\). При \(d=2\) добавляется блок \(D\) размера 2 на глубине 1. При \(d\ge3\) добавляется блок \(\left(\begin{smallmatrix}0&D\\D^*&0\end{smallmatrix}\right)\) на глубинах \(1,d-1\) и нуль на остальных \(2d-6\) координатах.
\newpage
\sect{8. Две сингулярные величины основного блока}
В блоке \(C\) подставим \(Y=az\) и \(q=ab=Q/(-2\rho)\), где \(Q=v_1^2-2v_2^2\). Тогда стационарное уравнение и критическая фаза равны
\[
  q z(7^d z-1)=\mu,\qquad
  \Phi_Q=q(-1/2+7^d z)+7^d\mu\big[\mathcal L(z)-\mathcal L(7^d z-1)\big].
                                                                  \tag{3}
\]
Стационарная точка единственна по Гензелю, \(z\equiv-\mu/q\bmod7\). Вторые производные фазы после снятия \(7^d\) принадлежат одному квадратному классу, так как по модулю 7 они равны \(-\mu/z^2\). Формула (2) поэтому даёт общий гауссов множитель \(\kappa\) для всех \(Q\), включая нормировку векторов, с \(|\kappa|=(7/3)N^{-1/2}\).

Обозначим через \(b_Q\) сумму коэффициентов \(c_v\) при данной норме \(Q\) и квадратном \(-a\bmod7\). При противоположном квадратном классе сумма равна \(\overline{b_Q}\), по переходу \(v\mapsto-v\). Вклад после снятия \(\kappa\) равен
\[
  \omega(\mu/q)^h\zeta_{7^k}^{\Phi_Q}
                      \diag(b_Q,\overline{b_Q})\,P^{(1-\epsilon_Q)/2},
  \qquad \epsilon_Q=\left(\frac{\mu/q}{7}\right).
\]
Здесь \(P\) меняет два столбца, а \((\cdot/7)\) --- квадратичный характер. Действительно, строка определяется квадратным классом \(x\equiv-a\), столбец классом \(Y\equiv az\), их отношение по модулю 7 равно \(\mu/q\). Все \(11\) сумм \(b_Q\) ненулевые и имеют рациональные вещественную и мнимую части.

Положим \(C'=C/\kappa\). Во внедиагональный элемент \(C'C'^*\) входят упорядоченные пары \((Q,R)\) с \(\epsilon_Q\ne\epsilon_R\). Их коэффициенты ненулевые, а фазы равны \(\Phi_Q-\Phi_R\). В \(\det C'\) входят неупорядоченные пары с одинаковой \(\epsilon\) и \(\Re(b_Q\overline b_R)\ne0\), включая \(Q=R\). Их фазы равны \(\Phi_Q+\Phi_R\), коэффициенты содержат соответственно \(2\Re(b_Q\overline b_R)\) или \(|b_Q|^2\) с фиксированным знаком и ненулевым весом \(\omega\).

\subt{Сертификат фаз}
По модулю \(343\) код находит по одному показателю кратности 1 в каждой из этих двух сумм. При \(d=1\) для каждого \(\mu\bmod343\) используется точная формула (3). При \(d\ge2\), \(\Phi_Q\equiv-q/2-7^d\mu\bmod343\), и общий сдвиг не влияет на кратности. В обоих случаях перечислено по \(294\) обратимых \(\mu\). Полная таблица фаз и выбранных пар включена в архив. Циклотомическая независимость даёт \(\det C'\ne0\) и \((C'C'^*)_{+-}\ne0\). Поэтому \(C\) имеет две положительные, различные сингулярные величины; его эрмитов блок обладает четырьмя простыми ненулевыми собственными числами.
\newpage
\sect{9. Отделение остальных блоков}
При \(d\ge3\) блок \(D\) между глубинами \(1,d-1\) получает только две непримитивные метки. После снятия постоянной части фазы и множителя \(7^d\) остаётся одно невырожденное стационарное слагаемое на метку, с условием того же вида, что (3), где \(7^d\) внутри логарифма заменяется на \(7^{d-2}\) и используются \(a_0,b_0\). В сумме имеется дополнительный множитель 7, поскольку \(Y\) пробегает \(7^{k-d+1}\) классов, а фаза от \(Y\) имеет период \(7^{k-d}\). Нормирующий знаменатель \((3/7)7^{k-d/2}\) такой же, как для \(C\). Две метки имеют коэффициент \(9/20\), противоположные квадратные классы строк и столбцов, поэтому обе сингулярные величины \(D\) равны \(7(9/20)|\kappa|\). С другой стороны,
\[
                  \|C\|\le |\kappa|\sum_{v\ {\rm примитивная}}|c_v|
                    <0.3|\kappa|.
\]
Здесь полная сумма модулей не больше \(593/500\), а две непримитивные метки дают \(9/10\). Спектр блока \(D,D^*\) отделён от спектра блока \(C,C^*\) и от нуля.

При \(d=2\) блок \(D\) действует на одной глубине, коммутирует с чётностью и самосопряжён. Поэтому имеет вид \(\alpha I+\beta P\), где \(\alpha,\beta\) лежат в
\(K'=\QQ(i,\zeta_3,\zeta_{7^{k-2}})\): вся вейлева и логарифмическая фаза делится на \(7^2\), а нормировка в сжатии рациональна. Покажем, что ни одно собственное число \(\lambda\) этого блока не может входить в спектр основного блока. Иначе \(u=\lambda^2/|\kappa|^2\in K'\) и
\[
                u^2-u\Tr(C'C'^*)+|\det C'|^2=0 .                 \tag{4}
\]

Фазы \(\Tr(C'C'^*)\) имеют вид \(\Phi_Q-\Phi_R\) с \(\epsilon_Q=\epsilon_R\). Фазы последнего слагаемого (4) имеют вид \(\Phi_Q+\Phi_R-\Phi_U-\Phi_V\), по двум допустимым детерминантным парам. По модулю 343 во всех этих разностях фазы равны соответственно \(Q-R\) и \(Q+R-U-V\), умноженным на одну и ту же единицу \(-1/[2(-2\rho)]\). Целочисленная проверка для каждого \(\mu\bmod7\) показывает, что пары \((-162,-162),(-97,1)\) дают единственный коэффициент при классе \(Q+R-U-V=115\bmod343\), а модуль \(49\) этого класса не встречается в следе. Он также ненулевой по модулю 49. Коэффициент при этом показателе в \(|\det C'|^2\) ненулевой по независимости степеней над \(\QQ(i,\zeta_3,\zeta_{7^{k-3}})\). Над \(K'\) умножение на \(u\) не меняет классы показателей по модулю 49. Следовательно, в (4) найденный класс не может обнулиться. Это доказывает отделение четырёх простых ветвей при \(d=2\) и завершает расчёт при каждом фиксированном \(d\ge1\).
\newpage
\sect{10. Следствия для фиксированного кика}
Доля кластеров порядка \(d\) на размерность равна \(18/7^{d+2}\) для каждого фиксированного \(d\). При \(d=0\) мы получаем \(249/392+O(N^{-1/2})\) простых уровней на размерность; при \(d\ge1\) в каждом кластере получаем не меньше четырёх. Сначала суммируем до произвольного фиксированного \(d_{\max}\), переходим к \(k\to\infty\), затем к \(d_{\max}\to\infty\):
\[
 \liminf \frac{s_N}{N}\ge\frac{249}{392}+\frac{72}{49}\sum_{d\ge1}7^{-d}
                         =\frac{345}{392}.
\]
Применён алгебраический перенос из раздела 1. Доказано существование всех указанных простых ветвей одновременно при достаточно малом ненулевом \(t\) для каждого отдельного \(N\), после чего тот же нижний предел справедлив при точном \(t=N/10^{80}\).

Первоначальная сумма квадратов кратностей есть число вейлевых циклов, то есть \(1+\sum_{a=1}^k(7^{2a}-7^{2a-2})/(3\cdot 7^{a-1})=(8N-5)/3\). Для кластеров порядка \(d=0\) расщепление двух ветвей уменьшает её на 2 в каждом расщеплённом кластере. Для положительных \(d\) уменьшение не меньше
\[
 \begin{cases}
 12,&d=1,\\
 28,&d=2,\\
 32d-44,&d\ge3.
 \end{cases}
\]
Действительно, при \(d=1\) кластер размера 4 разделяется на четыре простых, при \(d=2\) остаются четыре простых и не более одной пары. При \(d\ge3\) остаются четыре простых, две пары и блок нулевых производных размера \(2d-6\). Линейные кластеры при малом \(t\) разделены, так что остальные кластеры сумму квадратов не увеличивают. При фиксированном алгебраическом \(t\) ранг \(\Ad(U_N(t))-I\) не меньше ранга при малом \(t\). Подстановка в (1) и суммирование до фиксированного \(d_{\max}\) с последующим \(d_{\max}\to\infty\) дают
\[
 \limsup_{k\to\infty}\lim_{T\to\infty}
     \frac1{T7^k}\sum_{m=1}^T|\Tr U_{7^k}^m|^2
 \le \frac83-\frac{249}{392}
  -\frac{18}{{49}}\left(\frac{12}{7}+\frac{28}{7^2}
                    +\sum_{d\ge3}\frac{32d-44}{7^d}\right)
 =\frac{64549}{57624}.
\]
Эта граница относится к последовательному пределу, где сначала бесконечное среднее по времени. Для парного теста BGS-F нужен иной режим: интервалы времён длиной порядка \(N\). Нынешний довод классифицирует точные совпадения частот. Количественное суммирование при ненулевых щелях длины порядка \(1/N\) остаётся открытым.
\newpage
\sect{11. Точная кривизна обычного следа}

Пусть \(D=\sum_{j=0}^{m-1}M_N^j A_N M_N^{-j}\). Из \(M_N^m=\eta I\), \(|\eta|=1\), и упорядоченного произведения экспонент точно следуют
\[
 \left.\frac{d}{dt}\frac{\overline\eta}{N}\Tr U_N(t)^m\right|_{0}
       =i\Tr D/N=0,\qquad
 \left.\frac{d^2}{dt^2}\frac{\overline\eta}{N}\Tr U_N(t)^m\right|_{0}
       =-\Tr D^2/N .
\]
Во втором тождестве цикличность следа объединяет два порядка пар. Для \(k\ge3\) все исходные метки ненулевые, так что \(\Tr A_N=0\).

Орбиты примитивных меток характеризуются целой нормой \(Q(v)=v_1^2-2v_2^2\) и квадратным классом \(v_1+\rho v_2\bmod7\). Действительно, равенство норм по модулю \(N\) эквивалентно их целому равенству для этого носителя и \(N\ge343\). В расщеплённых координатах вторая координата определяется первой и нормой, а допустимые отношения первых координат образуют \(\langle\gamma\rangle\), ровно единицы с квадратным первым разрядом. Две непримитивные метки делим на \(7\) и проверяем то же сравнение по модулю \(N/7\); они имеют разные квадратные классы. Длины орбит \(L\) равны \(m\) или \(m/7\). В \(D\) коэффициент каждой выходной метки данной орбиты равен \((m/L)\sum_{v\ {\rm в\ орбите}}c_v\). По ортогональности Вейля
\[
   N^{-1}\Tr D^2
       = \sum_{\mathcal C}\frac{m^2}{L_{\mathcal C}}
                           \left|\sum_{v\in\mathcal C}c_v\right|^2
       =m\cdot\frac{11370883}{4000000}.
\]
Последнее равенство --- точная группировка исходных 32 коэффициентов в 24 класса; код сохраняет каждую сумму в гауссовых целых после умножения на \(4000\). Нормированный квадрат следа имеет вторую производную по \(t\), равную \(-2N^{-1}\Tr D^2\). Умножение на \(N^2\) при переходе к физической силе \(u=t/N\) даёт заявленную формулу. Точное расщепление и кривизна не дают оценки остатка при \(t=N/10^{80}\) и \(N\to\infty\).

\subt{Проверяемость}
В архив включены таблица стационарных свидетельств, независимая циклотомическая проверка малых \(k\), исходные коэффициенты и сборка документа. Команда \texttt{python3 spectrum/check.py --verify} проверяет все конечные утверждения целыми и рациональными числами.
\end{document}
